% SPDX-License-Identifier: CC-BY-4.0
\documentclass[11pt,letterpaper]{article}

\usepackage[T1]{fontenc}
\usepackage{amsmath,amssymb}
\usepackage{tikz}
\usepackage{pgfplots}
\pgfplotsset{compat=1.18}
\usepackage{algorithm}
\usepackage{algpseudocode}
\usepackage{booktabs}
\usepackage{enumitem}
\usepackage[hidelinks]{hyperref}

\newcommand{\N}{\mathbb{N}}
\newcommand{\Z}{\mathbb{Z}}
\newcommand{\Q}{\mathbb{Q}}
\newcommand{\R}{\mathbb{R}}
\newcommand{\C}{\mathbb{C}}
\newcommand{\Map}[3]{#1 \colon #2 \to #3}
\newcommand{\Mod}[1]{#1\text{-}\mathrm{Mod}}
\newcommand{\GroupHom}{\mathrm{Hom}_{\mathbf{Grp}}}
\newcommand{\RingHom}{\mathrm{Hom}_{\mathbf{Ring}}}

\newtheorem{theorem}{Theorem}
\newtheorem{lemma}[theorem]{Lemma}
\newtheorem{corollary}[theorem]{Corollary}
\newtheorem{definition}[theorem]{Definition}
\newtheorem{remark}[theorem]{Remark}
\newenvironment{proof}{\noindent\textit{Proof.}}{\hfill\(\square\)}

\title{Generalized Collatz Library}
\author{}
\date{\today}

\begin{document}
\maketitle
\section*{Three-Affine Closure and Exact Odd Recovery}

\begin{abstract}
For positive integers \(a,b,c,q\) and integers \(x,y,z\) linked by one pass of the three affine letters \(2^{a}y=qx+1\), \(2^{b}z=qy+1\), and \(2^{c}x=qz+1\), the determinant \(\Delta=2^{a+b+c}-q^{3}\) divides the closure numerator \(B=q^{2}+q2^{a}+2^{a+b}\). The third odd state is the same quotient formed from the rotated numerator \(q^{2}+q2^{c}+2^{c+a}\). For positive odd states, we also prove the converse: a legal closure exists exactly when the multiplier is odd, the determinant is positive, and the determinant divides the closure numerator. Rotated numerators reconstruct all phases and their exact powers of two.
\end{abstract}

\section{Closure numerator}

Let \(a,b,c,q\) be positive integers and set \(K=a+b+c\). Write
\[
B(a,b;q)=q^{2}+q\,2^{a}+2^{a+b},
\qquad
T(c,a;q)=q^{2}+q\,2^{c}+2^{c+a},
\]
and \(\Delta(a,b,c;q)=2^{K}-q^{3}\). Thus \(B(a,b;q)\) is the closure numerator of the word \((a,b,c)\), and \(T(c,a;q)\) is the closure numerator of the rotated word \((c,a,b)\).

\begin{lemma}
\label{lem:numerator-seven}
If \(e,f,r\) are positive integers, then
\[
r^{2}+r\,2^{e}+2^{e+f}\ge 7.
\]
In particular \(B(a,b;q)\ge 7\) and \(T(c,a;q)\ge 7\).
\end{lemma}

\begin{proof}
Since \(r\ge 1\), one has \(r^{2}\ge 1\). Since \(e\ge 1\), one has \(2^{e}\ge 2\), and therefore \(r\,2^{e}\ge 2\). Since \(e+f\ge 2\), one has \(2^{e+f}=4\cdot 2^{e+f-2}\ge 4\). Adding these three estimates gives
\[
r^{2}+r\,2^{e}+2^{e+f}\ge 1+2+4=7.
\]
The two particular cases are \((e,f,r)=(a,b,q)\) and \((e,f,r)=(c,a,q)\).
\end{proof}

\begin{lemma}
\label{lem:identity}
Let \(e_{1},e_{2},e_{3}\) be natural numbers and let \(q,u,v,w\in\Z\) satisfy
\[
2^{e_{1}}u=qv+1,
\qquad
2^{e_{2}}v=qw+1,
\qquad
2^{e_{3}}w=qu+1.
\]
Then
\[
\bigl(2^{e_{1}+e_{2}+e_{3}}-q^{3}\bigr)u
  =q^{2}+q\,2^{e_{3}}+2^{e_{3}+e_{2}}.
\]
\end{lemma}

\begin{proof}
Multiply the first equation by \(2^{e_{2}+e_{3}}\):
\[
2^{e_{1}+e_{2}+e_{3}}u
  =2^{e_{2}+e_{3}}(qv+1)
  =q\,2^{e_{2}+e_{3}}v+2^{e_{2}+e_{3}}.
\]
Split the remaining power and apply the second equation:
\[
2^{e_{2}+e_{3}}v=2^{e_{3}}(qw+1),
\]
so
\[
2^{e_{1}+e_{2}+e_{3}}u
  =q\,2^{e_{3}}(qw+1)+2^{e_{2}+e_{3}}
  =q\bigl(q\,2^{e_{3}}w+2^{e_{3}}\bigr)+2^{e_{2}+e_{3}}.
\]
The third equation replaces \(2^{e_{3}}w\) by \(qu+1\):
\begin{align*}
2^{e_{1}+e_{2}+e_{3}}u
  &=q\bigl(q(qu+1)+2^{e_{3}}\bigr)+2^{e_{2}+e_{3}}\\
  &=q^{3}u+q^{2}+q\,2^{e_{3}}+2^{e_{2}+e_{3}}.
\end{align*}
Subtract \(q^{3}u\). The exponent \(e_{2}+e_{3}\) may be written \(e_{3}+e_{2}\), which is the stated numerator.
\end{proof}

\begin{theorem}
\label{thm:three-affine-closure}
Let \(a,b,c,q\) be positive integers and let \(x,y,z\in\Z\) satisfy
\[
2^{a}y=qx+1,
\qquad
2^{b}z=qy+1,
\qquad
2^{c}x=qz+1.
\]
Set \(\Delta=\Delta(a,b,c;q)\), \(B=B(a,b;q)\), and \(T=T(c,a;q)\). Then \(\Delta x=B\), \(\Delta\) divides \(B\), \(\Delta z=T\), \(z=T/\Delta\), and \(\Delta\neq 0\). Both numerators are at least \(7\).
\end{theorem}

\begin{proof}
Apply Lemma \ref{lem:identity} to the ordered equations
\[
2^{c}x=qz+1,
\qquad
2^{b}z=qy+1,
\qquad
2^{a}y=qx+1.
\]
The resulting coefficient is \(2^{c+b+a}-q^{3}=\Delta\), and the resulting numerator is \(q^{2}+q\,2^{a}+2^{a+b}=B\). Thus \(\Delta x=B\), so \(\Delta\) divides \(B\).

Apply the same lemma to
\[
2^{b}z=qy+1,
\qquad
2^{a}y=qx+1,
\qquad
2^{c}x=qz+1.
\]
The coefficient is again \(\Delta\), and the numerator is \(q^{2}+q\,2^{c}+2^{c+a}=T\). Thus \(\Delta z=T\).

Lemma \ref{lem:numerator-seven} gives \(B\ge 7\). If \(\Delta=0\), then \(B=\Delta x=0\), contradicting \(B\ge 7\). Hence \(\Delta\neq 0\), and dividing \(\Delta z=T\) by \(\Delta\) yields \(z=T/\Delta\). The same lower bound gives \(T\ge 7\).
\end{proof}

\section{Positive odd closure}

Put \(U=q^2+q2^b+2^{b+c}\) and \(V=q^2+q2^c+2^{c+a}\).
Expanding both sides gives
\[
qB+\Delta=2^aU,\qquad
qU+\Delta=2^bV,\qquad
qV+\Delta=2^cB.
\]
For example,
\[
qB+\Delta=q^2 2^a+q2^{a+b}+2^{a+b+c}
=2^a(q^2+q2^b+2^{b+c}).
\]
The remaining two identities follow by cyclic permutation.

% Rotated divisibility recovers all three positive odd phases, not just the first.
\begin{theorem}
Let \(a,b,c,q\) be positive integers. Positive odd integers \(x,y,z\)
satisfying the three equations exist if and only if
\[
q\text{ is odd},\qquad \Delta>0,\qquad \Delta\mid B.
\]
When they exist, the ordered triple is uniquely
\[
(x,y,z)=(B/\Delta,U/\Delta,V/\Delta).
\]
\end{theorem}

\begin{proof}
For necessity, \(2^a y\) is even and \(x\) is odd. Thus \(qx+1\)
can be even only when \(q\) is odd. The elimination identity gives
\(\Delta x=B\). Since \(x>0\) and \(B>0\), it follows that
\(\Delta>0\) and \(\Delta\mid B\).

For sufficiency, an odd \(q\) makes \(\Delta\) odd and nonzero.
Hence \(\gcd(\Delta,2^a)=\gcd(\Delta,2^b)=1\).
The first transport identity and \(\Delta\mid B\) give
\(\Delta\mid 2^aU\), so \(\Delta\mid U\).
The second identity then gives \(\Delta\mid V\).
Define \(x=B/\Delta\), \(y=U/\Delta\), and \(z=V/\Delta\).
All numerators are positive. Each numerator is odd because its first
term is odd and its other two terms are even. Dividing by a positive odd
divisor gives positive odd integer quotients.
Dividing the three transport identities by \(\Delta\) now proves all
three required equations. Their odd quotient states give
\(v_2(qx+1)=a\), \(v_2(qy+1)=b\), and \(v_2(qz+1)=c\).
Elimination at each phase forces the displayed quotients, proving uniqueness.
\end{proof}

Each recovered odd-return edge consists of one affine raw tick followed
by exactly its prescribed number of divisions. The three-edge return
therefore has raw length \(K+3\) and shortcut length \(K\).
These return lengths need not be primitive periods.

For \(q=5\) and \((a,b,c)=(1,1,5)\), one obtains
\[
\Delta=3,\quad (B,U,V)=(39,99,249),\quad (x,y,z)=(13,33,83).
\]
The odd-return cycle is \(13\to33\to83\to13\).
The parity condition cannot be omitted: for \(q=2\) and
\((a,b,c)=(1,2,1)\), one has \(\Delta=8\mid B=16\), but the first
equation forces \(y=5/2\).

The theorem is an exact criterion for each fixed positive parameter
word. It does not impose a uniform upper bound on \(K\), and does not
classify every parameter word that passes the arithmetic tests.
The elimination identities are formalized in \texttt{ThreeAffineClosure.lean};
the positive odd criterion and rotated recovery are formalized in
\texttt{ThreeAffineLegal.lean}.

\section{Classical short return closures}

The general criterion leaves arbitrarily large exponents. For the classical
multiplier, the product of the edge equations supplies a bounded exponent
window when the number of odd-return edges is at most five.

% The statement counts return edges and allows repetition of the unit cycle.
\begin{theorem}[At most five classical edges]
Let \(1\le r\le5\). Suppose positive odd integers \(x_0,\ldots,x_{r-1}\)
and positive integers \(a_0,\ldots,a_{r-1}\) satisfy
\[
2^{a_i}x_{i+1}=3x_i+1
\quad(0\le i<r),
\]
with cyclic indices. Then every \(x_i=1\) and every \(a_i=2\).
\end{theorem}
\begin{proof}
If any \(x_i=1\), the next edge has \(2^{a_i}x_{i+1}=4\).
Its positive odd quotient forces \(a_i=2\) and \(x_{i+1}=1\).
Propagation around the finite cycle proves the assertion.

Otherwise every state is at least three. Multiply all edge equations
and cancel the positive product of the states. With \(S=\sum_i a_i\),
\[
3^r<2^S=\prod_{i=0}^{r-1}\left(3+\frac1{x_i}\right)
\le\left(\frac{10}{3}\right)^r.
\]
For \(r=1,2,4\) this interval contains no power of two: its endpoints
lie respectively between \(2\) and \(4\), between \(8\) and \(16\),
and between \(64\) and \(128\). The explicit windows are
\((3,10/3]\), \((9,100/9]\), and \((81,10000/81]\).

For \(r=3\), the interval forces \(S=5\), since
\(27<2^S\le1000/27<64\). The elimination identity gives
\[
5x_0=9+3\cdot2^{a_0}+2^{a_0+a_1}.
\]
The six positive compositions of five into three parts give numerators
\[
19,23,31,29,37,49.
\]
None is divisible by five, a contradiction.

For \(r=5\), the window forces \(S=8\), since
\(243<2^S\le100000/243<512\). Successive substitution gives
\[
13x_0=N(a),\qquad
N(a)=\sum_{i=0}^4 3^{4-i}2^{a_0+\cdots+a_{i-1}},
\]
where the empty exponent sum is zero. The 35 positive compositions of
eight into five parts have the following seven cyclic representatives.
\begin{center}
\begin{tabular}{lrr}
\toprule
Exponents & \(N(a)\) & \(N(a)\bmod13\) \\
\midrule
\((1,1,1,1,4)\) & 211 & 3 \\
\((1,1,1,2,3)\) & 227 & 6 \\
\((1,1,1,3,2)\) & 259 & 12 \\
\((1,1,2,1,3)\) & 251 & 4 \\
\((1,1,2,2,2)\) & 283 & 10 \\
\((1,1,3,1,2)\) & 331 & 6 \\
\((1,2,1,2,2)\) & 319 & 7 \\
\bottomrule
\end{tabular}
\end{center}
To check exhaustiveness, subtract one from each exponent. The remaining
sum is three, with types \(3\), \(2+1\), and \(1+1+1\).
They have respectively one, four, and two cyclic arrangements: the two
zeros in the last type are either adjacent or separated on the five-cycle.
Each arrangement has five rotations, because a nonconstant five-letter
word has no shorter cyclic period. Thus the table covers all 35 cases.
A hypothetical cycle can be rotated to the corresponding row, retaining
its legal equations. No row has the required divisibility by thirteen.
\end{proof}

This classifies closures with up to five odd-return edges, not arbitrary
positive cycles. Their shortcut return length is \(\sum_i a_i\); the
raw length is \(\sum_i a_i+r\). The conclusion has shortcut length \(2r\)
and raw length \(3r\), repeating the primitive unit cycle.
The companion
\nolinkurl{QCD.ThreeAffineClosure.001/lean/ClassicalThreeReturn.lean}
formalizes the complete three-edge classification. It also formalizes
five-edge elimination and the failure of all 35 positive exponent-eight
compositions. The five-edge product window and the full at-most-five
statement above are written proofs.

\end{document}
